% 01-introduction
\chapter{Introduction: patterns and Turing's idea}\label{ch:introduction}

Many structures in nature are neither random nor uniform: they repeat some
element---a spot, a stripe, a hair, a finger---at a fairly regular distance. In
1952 Alan Turing proposed that such patterns can arise from chemical reactions and
diffusion alone. This chapter describes the phenomenon and the mechanism, and asks
the questions that these notes answer.

\section{Patterns everywhere}\label{sec:intro-patterns}

Think of the spots of a leopard, the stripes of a zebra, or the stripes and spots
on the skin of many fish, which in some species even rearrange as the animal grows.
Regularly spaced structures also appear in embryonic development: hair follicles and
feather buds are laid out at nearly equal distances, and a limb forms a row of
digits (see the survey \cite{KondoMiura2010}). The cleanest evidence comes from
chemistry. In 1990
Castets, Dulos, Boissonade and De Kepper ran the chlorite--iodide--malonic acid
(CIMA) reaction in a thin layer of gel continuously fed with fresh reactants, and
observed stationary spots and stripes whose spacing was chosen by the chemistry
itself \cite{Castets1990}. Lengyel and Epstein derived a two-variable model of this
reaction, in which binding in the gel makes the two key species diffuse at very
different effective rates \cite{LengyelEpstein1991}; it is one of our running
examples, \eqref{eq:lengyel-epstein}.

In all these examples the pattern seems to come out of nothing. The early tissue or
the fresh gel is nearly uniform, the laws that govern it treat all points alike, and
yet the final state is not uniform: this is \emph{symmetry breaking}. Moreover, the
distance between neighboring spots or stripes is roughly fixed, an \emph{intrinsic
length scale}. A theory should explain both.

\section{Turing's idea}\label{sec:intro-turing}

In \emph{The chemical basis of morphogenesis} \cite{Turing1952}, Turing suggested
that a few substances, which he called \emph{morphogens}, ``reacting together and
diffusing through a tissue'' could account for the formation of biological
structure. This contains a paradox. \emph{Diffusion alone} is the great equalizer:
a solution of the diffusion equation with zero-flux boundary conditions converges to
its spatial mean (\cref{thm:heat-decay}). \emph{Reactions alone} are often
stabilizing, too: if a well-stirred chemical system has a stable equilibrium, the
concentrations return to it after a small disturbance. One would expect the
combination to be even more stable. Turing's discovery is that this intuition is
wrong: if two species diffuse at sufficiently different rates, diffusion can make the
uniform equilibrium \emph{unstable}, and the instability prefers certain
wavelengths. This is called \emph{diffusion-driven instability (DDI)} or
\emph{Turing instability} (\cref{def:ddi}).

\section{Local activation, long-range inhibition}\label{sec:intro-activation}

\begin{figure}[t]
  \centering
  \begin{tikzpicture}[
      species/.style={circle, very thick, minimum size=14mm, align=center},
      lab/.style={font=\footnotesize, text=black!75}]
    \node[species, draw=keycolor, fill=keycolor!8] (u) at (0,0)
      {\color{keycolor}$u$\\[-3pt]\color{keycolor}\scriptsize activator};
    \node[species, draw=red!70!black, fill=red!6] (v) at (5,0)
      {\color{red!70!black}$v$\\[-3pt]\color{red!70!black}\scriptsize inhibitor};
    \draw[-{Stealth[length=2.4mm]}, thick, keycolor]
      (u.150) .. controls +(150:1.3) and +(210:1.3) .. (u.210)
      node[lab, pos=0.5, left] {self-enhancement};
    \draw[-{Stealth[length=2.4mm]}, thick, keycolor]
      (u.25) to[bend left=18] node[lab, above] {activates} (v.155);
    \draw[-{Bar[width=3mm]}, thick, red!70!black]
      (v.205) to[bend left=18] node[lab, below] {inhibits} (u.335);
    \draw[{Stealth[length=2mm]}-{Stealth[length=2mm]}, thick, keycolor]
      (0.5,1.0) -- (-0.5,1.0) node[lab, left] {slow diffusion};
    \draw[{Stealth[length=2mm]}-{Stealth[length=2mm]}, thick, red!70!black]
      (3.6,1.0) -- (6.4,1.0) node[lab, right] {fast diffusion};
  \end{tikzpicture}
  \caption{Local activation and long-range inhibition: the activator $u$ enhances
  itself and the inhibitor $v$, which suppresses $u$. The double arrows indicate the
  range of diffusion.}
  \label{fig:ch01-activator-inhibitor}
\end{figure}

Why should diffusion destabilize? The best-known explanation is due to Gierer and
Meinhardt \cite{GiererMeinhardt1972} (\cref{fig:ch01-activator-inhibitor}). An
\emph{activator} $u$ promotes its own production and that of an \emph{inhibitor}
$v$, which suppresses $u$. Suppose the activator is slightly elevated somewhere.
Self-enhancement makes the excess grow, and more inhibitor is produced there as
well. If the inhibitor stayed where it was produced, it would quench the peak. If
it diffuses much faster than the activator, it leaks into the
surroundings: at the peak the activator wins, around it the inhibitor keeps the
activator low, and farther away the next peak can form. The result is a periodic
arrangement of peaks whose spacing is set by the range of the inhibitor. A
fast-diffusing \emph{substrate} that the activator consumes, and hence depletes
around a peak, has the same effect (\cref{ch:ode-toolbox} has both variants).

Think of plants in a dry landscape: they help each other locally by shading the
ground, but their roots compete for water over a wide area, and clumps form.

\begin{keyidea}
\emph{Local activation and long-range inhibition}: a slowly diffusing,
self-enhancing species coupled to a fast-diffusing antagonist can turn a uniform
state into a regular pattern. \Cref{ch:turing} turns this slogan into theorems
(\cref{thm:ddi-criterion,cor:turing-consequences}).
\end{keyidea}

\section{A first look at the equations}\label{sec:intro-model}

For the concentrations $u(x,t)$, $v(x,t)$ of two species we study models of the
form
\begin{equation*}
  u_t = d_u\,u_{xx} + f(u,v),\qquad
  v_t = d_v\,v_{xx} + g(u,v),\qquad
  x\in(0,L),\ t>0,
\end{equation*}
with $u_x=v_x=0$ at $x=0,L$: nothing leaves or enters. The concentrations change
because molecules spread by diffusion, with diffusion coefficients $d_u,d_v>0$
(\cref{ch:diffusion}), and are produced or destroyed by reactions, described by $f$
and $g$. With initial data this is the system \eqref{eq:rd-system} of
\cref{ch:turing}. Constants $(\ub,\vb)$ with $f(\ub,\vb)=g(\ub,\vb)=0$ solve it:
a uniform state.

\begin{keyidea}[Three questions]
\begin{enumerate}[label=(Q\arabic*), leftmargin=3.2em]
  \item When is a constant steady state stable without diffusion but unstable with it?
  \item Which spatial wavelengths grow, and what is the role of the domain size $L$?
  \item What happens after the instability? Does a stationary pattern form?
\end{enumerate}
\end{keyidea}

(Q1) is answered in \cref{ch:one-equation,ch:turing}: one equation is never
enough; for two, the answer is the Turing conditions \eqref{eq:turing-conditions}.
(Q2) is answered in \cref{ch:turing} (\cref{cor:domain-size}) and illustrated in
\cref{ch:examples}. (Q3) is nonlinear and answered only in part, by simulations in
\cref{ch:examples} and a survey in \cref{ch:outlook}.

\section{Road map and how to read these notes}\label{sec:roadmap}

\Cref{ch:ode-toolbox} collects tools from ordinary differential equations and
introduces the running examples (Schnakenberg, Gierer--Meinhardt and
Lengyel--Epstein kinetics). \Cref{ch:diffusion} treats diffusion alone and solves
the diffusion equation with the Neumann eigenfunctions $\cos(k\pi x/L)$.
\Cref{ch:one-equation} shows that one reaction--diffusion equation cannot exhibit
diffusion-driven instability and that its nonconstant stationary solutions are
unstable. \Cref{ch:turing} is the core: the exact criterion for diffusion-driven
instability with two species, the critical diffusion ratio, and the role of the
domain size. \Cref{ch:examples} maps the parameter regions where the running
examples form patterns (Turing spaces) and simulates the nonlinear equations.
\Cref{ch:outlook} explains what lies beyond. \Cref{app:fourier,app:linear-algebra}
collect facts on cosine series and on $2\times2$ matrices.

We assume a course on ordinary differential equations (linear systems, matrix
exponential, linearization, phase planes), real analysis and linear algebra, but no
knowledge of partial differential equations. We always work on an interval with
homogeneous Neumann (zero-flux) boundary conditions, natural for a closed system.
The eigenfunctions of the second derivative are then explicit cosines, so Fourier
cosine series replace the spectral theory of the Laplace operator on a general
domain; the results remain true on bounded domains in $\R^n$
(\cref{rem:general-domain}). Key-idea boxes summarize the central take-aways,
exercises marked ($\ast$) are harder, and all figures are made by short Python
scripts in \texttt{scripts/}, ready to be run and modified.

\section{Exercises}\label{sec:intro-exercises}

\begin{exercise}[Units and time scales]\label{exer:intro-units}
(a) Show that in $u_t=d\,u_{xx}$ the coefficient $d$ has the unit
$\text{length}^2/\text{time}$, so that $\ell^2/d$ is a time for every length $\ell$.
(b) Assume $d\approx10^{-11}\,\mathrm{m^2/s}$ for a protein in tissue. Estimate
$\ell^2/d$ for $\ell=0.1\,\mathrm{mm}$, $1\,\mathrm{mm}$ and $1\,\mathrm{cm}$. What
does this suggest about the size of tissues in which such molecules lay down a
pattern?
\end{exercise}

\begin{exercise}[Why not one species?]\label{exer:intro-one-species}
Let $u_t=d\,u_{xx}+f(u)$ on $(0,L)$, $u_x=0$ at $x=0,L$, and let $f(\ub)=0$,
$f'(\ub)<0$. Show that $w(x,t)=e^{\mu t}\cos(k\pi x/L)$ solves the linearized
equation $w_t=d\,w_{xx}+f'(\ub)\,w$ with the same boundary condition if and only if
$\mu=f'(\ub)-d\,(k\pi/L)^2$. Why does this suggest that one species cannot produce a
Turing instability? (See \cref{cor:no-ddi-scalar,thm:nonconstant-unstable}.)
\end{exercise}

\begin{exercise}[Reading an interaction network]\label{exer:intro-signs}
Determine the signs of $f_u,f_v,g_u,g_v$ at a steady state, assuming that
self-enhancement is stronger than degradation:
(a) $u$ stimulates its own production and that of $v$, $v$ represses the production
of $u$, and both decay; (b) a substrate $v$ is supplied at a constant rate and
converted into $u$ in a reaction catalyzed by $u$, and $u$ decays.
In each case, which species should diffuse faster for a pattern to form? Compare
with \eqref{eq:gierer-meinhardt} and \eqref{eq:schnakenberg}.
\end{exercise}
